% Why Are HPC Monitoring Signals Hard to Forecast?
% Working draft — arXiv format, venue not yet chosen.
%
% Build: make paper   (pdflatex + bibtex; latexmk is not installed on the
% development machine, so the Makefile chains the passes explicitly)
%
% Every number in this document must be traceable to a committed table under
% results/tables/ or data/derived/. Placeholders are marked \needsresult so
% they are impossible to miss in the compiled PDF.

\documentclass[11pt,a4paper]{article}

\usepackage[utf8]{inputenc}
\usepackage[T1]{fontenc}
\usepackage{lmodern}
\usepackage[margin=1in]{geometry}
\usepackage{amsmath,amssymb}
\usepackage{booktabs}
\usepackage{graphicx}
\usepackage{xcolor}
\usepackage[hidelinks]{hyperref}
\usepackage{natbib}

% Number and unit helpers.
%
% siunitx is deliberately NOT used. It is absent from a minimal TeX Live
% install, and a paper that will not compile without an extra package is a
% reproducibility problem for exactly the readers most likely to try. These do
% the little this document actually needs.
\newcommand{\num}[1]{#1}
\newcommand{\unit}[2]{#1\,#2}
\newcommand{\unitrange}[3]{#1--#2\,#3}

% Marks a value that has not been computed yet. Deliberately loud.
\newcommand{\needsresult}[1]{\textcolor{red}{\textbf{[#1]}}}
% Cites an entry in OPEN_ISSUES.md, keeping paper and register in sync.
\newcommand{\oi}[1]{\textsc{oi-#1}}

\title{Why Are HPC Monitoring Signals Hard to Forecast?\\
       \large Intrinsic Predictability Across Nine Telemetry Families
       of a Tier-0 Supercomputer}

\author{Mohsen Seyedkazemi Ardebili\\
        \small \texttt{mohsen.seyedkazemi@gmail.com}}

\date{Working draft --- \today}

\begin{document}
\maketitle

\begin{abstract}
Operational telemetry from high-performance computing systems resists
forecasting to a degree that is widely reported anecdotally and rarely
explained. The usual response is to increase model capacity. We argue instead
that the difficulty is largely \emph{intrinsic to the signals}, differs
systematically between signal families, and can be quantified without fitting
any forecaster.

Using the public M100 ExaData release from CINECA's Marconi100 Tier-0
supercomputer, we characterise telemetry from nine heterogeneous plugin
families --- out-of-band sensors, operating-system counters, scheduler
aggregates, alarm states, facility power, cooling-plant control registers, and
exogenous weather --- collected from the same machine over the same 31 months,
so that signal type varies while system, workload, and period are held fixed.
For each series we compute model-free predictability measures spanning entropy,
periodicity, long memory, non-stationarity, and quantisation, and separately
measure the forecast skill achieved by a tiered suite ranging from naive
baselines to gradient-boosted trees under a strictly causal rolling-origin
protocol.

We report three findings. First, the monitoring infrastructure is itself
non-stationary: of \num{573} distinct (plugin, metric) pairs observed across the
campaign, only \num{104} are present in every month, and those are exactly the
out-of-band sensor set. Second, \needsresult{predictability landscape result}.
Third, \needsresult{linking result: out-of-sample \(R^2\) under
leave-one-family-out cross-validation}.

\needsresult{Implications paragraph.}
\end{abstract}

\section{Introduction}
\label{sec:intro}

\needsresult{To be written.} Framing points to develop:

\begin{itemize}
  \item Operators increasingly collect holistic telemetry, and increasingly
        want to \emph{predict} it --- for thermal safety, power capping, and
        proactive maintenance --- not merely to detect anomalies after the fact.
  \item Reported forecast accuracy on this class of signal is poor relative to
        standard time-series benchmarks, and the literature treats this as a
        modelling deficiency.
  \item We invert the question: rather than asking which model does best, we
        ask how much predictability is present at all, measured independently
        of any model.
  \item Contributions: (i) a quantitative census of signal heterogeneity in a
        Tier-0 monitoring campaign, including the instability of the monitoring
        system itself; (ii) a model-free predictability profile per signal
        family; (iii) a causal, scale-free forecast benchmark on the same
        series; (iv) a test of whether the first predicts the second, validated
        by holding out entire signal families.
\end{itemize}

\section{Background}
\label{sec:background}

\subsection{Marconi100 and the ExaData campaign}

\needsresult{Machine description and campaign scope.}

\subsection{The nine plugin families}
\label{sec:families}

The value of this dataset for the present question is that it holds the system
fixed while varying the signal. Table~\ref{tab:families} summarises the
families and the mechanism that generates each.

\begin{table}[t]
\centering
\caption{Telemetry families and their generating mechanisms. Nominal cadences
are as documented by the dataset; measured cadences differ substantially for
some families (\S\ref{sec:census}).}
\label{tab:families}
\small
\begin{tabular}{llll}
\toprule
Family & Measures & Nominal cadence & Generating mechanism \\
\midrule
IPMI      & BMC sensors: temperature, power, fans & \unit{20}{s} & thermal diffusion, quantised \\
Ganglia   & OS and GPU counters                   & \unitrange{20}{300}{s}, irregular & software instrumentation \\
SLURM     & cluster scheduler aggregates          & \unitrange{5}{10}{s} & user submission behaviour \\
Nagios    & service and hardware alarm states     & \unit{15}{min} & discrete state machine \\
Job table & per-job submission records            & event-driven & not a time series (\oi{8}) \\
Schneider & liquid-cooling PLC registers          & \unit{20}{s} & closed-loop control \\
Logics    & room power, PUE, chillers             & --- & facility aggregation \\
Vertiv    & air-handling units                    & \unit{10}{s} & closed-loop control \\
Weather   & outdoor conditions                    & \unit{10}{min} & exogenous \\
\bottomrule
\end{tabular}
\end{table}

\section{Related work}
\label{sec:related}

\needsresult{To be written from the screening table in
\texttt{literature/screening.csv}.} Every claim in this section must be
verified against the source PDF before submission; none has been, and this is
tracked as \oi{10}.

\section{Method}
\label{sec:method}

\subsection{Campaign census}
\label{sec:census-method}

\needsresult{Describe the inventory and per-series census.}

\subsection{The entity panel}

Node-level series of one metric are highly redundant --- nodes share racks,
inlet air, cooling zones, and workload --- and are therefore not independent
samples (\oi{9}). Analysis runs on a fixed panel of \num{96} entities whose
membership is a stable hash of the entity label under the project seed. Because
the hash sees only the label, panel membership is identical across every metric
and every month, so cross-metric comparisons are made on the same physical
machines rather than on a fresh random slice.

\subsection{A-priori signal taxonomy}

Labels are assigned from the plugin and metric \emph{name} alone, before any
predictability is measured, so that a finding of between-family differences
cannot be a relabelling of the result it is meant to explain.

\subsection{Model-free predictability measures}

\needsresult{Table of measures and what each probes.}

\subsection{Forecast protocol}

Skill is measured under rolling-origin evaluation with strict causality
asserted per split. Errors are scaled by MASE with a denominator computed from
the training window only; scaling by test-window statistics is the most common
silent leak in forecasting benchmarks and would compress exactly the
between-tier differences this study reads. The seasonal period is fixed per
series \emph{a priori} from its own cadence, with a non-seasonal robustness
check (\oi{5}).

MASE is undefined where the in-sample seasonal-naive error vanishes. This is
not an edge case here: a substantial share of control-plane series are constant
(\S\ref{sec:census}), and near-perfectly periodic windows yield denominators at
the level of floating-point noise. Such series are reported as undefined rather
than assigned a spuriously large score.

\subsection{Linking analysis}

\needsresult{Describe leave-one-family-out CV, baseline, ablations.}

\section{Results}
\label{sec:results}

\subsection{The monitoring system is not stationary}
\label{sec:census}

Across the 31-month campaign we catalogue \num{10504} metric partitions
comprising \num{573} distinct (plugin, metric) pairs. Only \num{104} --- all of
them IPMI --- are present in every month. Ganglia's metric count runs
\num{0}, \num{33}, \num{73}, \num{49}, \num{157}, \num{177} and is not
monotonic; SLURM's runs \num{20}, \num{2}, \num{0}, \num{2}, \num{36};
Schneider first appears in 21-03 and Vertiv and Weather in 21-04.

Any claim compared across months therefore confounds signal behaviour with
changes to the instrument, and we restrict cross-month statements accordingly
(\oi{2}).

\needsresult{Figure: coverage over time.}

\subsection{Signal heterogeneity}

\needsresult{Census results: cadence, cadence irregularity, coverage,
quantisation, constant fraction, per family.}

\subsection{The predictability landscape}

\needsresult{Feature distributions per family.}

\subsection{Forecast skill}

\needsresult{Skill by tier, horizon, and family; the saturation test.}

\subsection{Does predictability explain skill?}

\needsresult{Leave-one-family-out \(R^2\), importances, ablations.}

\section{Threats to validity}
\label{sec:threats}

The full register is maintained in \texttt{OPEN\_ISSUES.md} and each item is
cited from the text above. The load-bearing ones are:

\begin{description}
  \item[\oi{1}] Ganglia's sampling irregularity confounds instrumentation with
        signal physics. Sampling regularity is therefore carried as a covariate
        and reported separately.
  \item[\oi{2}] The monitoring system changes over the campaign.
  \item[\oi{6}] No neural tier has been run; no claim of the form ``not even
        deep models help'' is licensed.
  \item[\oi{7}] The error floor is an estimate with a known bias direction, not
        an information-theoretic bound.
  \item[\oi{11}] Permutation entropy is biased downward by quantisation, and
        quantisation varies by family, so entropy is interpreted only alongside
        the measured quantisation step.
\end{description}

\section{Conclusion}
\needsresult{To be written.}

\section*{Data and code availability}

All code, the versioned literature query set, the screening table, and every
derived table are at
\url{https://github.com/MSKazemi/hpc-signal-forecastability}. The underlying
measurements are the public M100 ExaData release~\citep{borghesi2023m100}.
Results reproduce from the committed derived tables without downloading raw
data.

\bibliographystyle{plainnat}
\bibliography{refs}

\end{document}
